\documentclass[11pt,a4paper]{article}
\usepackage[margin=1in]{geometry}
\usepackage{booktabs}
\usepackage{array}
\usepackage{enumitem}
\usepackage{hyperref}
\usepackage{xcolor}
\hypersetup{colorlinks=true,linkcolor=blue!50!black,urlcolor=blue!50!black}
\setlist{itemsep=2pt,topsep=4pt}
\renewcommand{\arraystretch}{1.25}

\title{\textbf{PakAgriVision: Crop Disease Detector}\\[4pt]\large Project Proposal -- Track A: Application Development}
\author{
  Abdul Ahad (24L-3029) \quad M. Anas (24L-3004) \quad M. Babar Shahzad (24L-3047)\\[4pt]
  GitHub: \url{https://github.com/ErenYeagerJinko/PakAgriVision}
}
\date{\today}

\begin{document}
\maketitle

\section{Problem Statement}
Farmers, especially in developing regions, often lack access to agricultural experts who can quickly and accurately diagnose plant diseases. By the time a disease is identified through traditional means, it has frequently spread beyond the point of effective treatment, leading to significant crop loss and economic hardship. There is a clear need for an accessible, instant, and trustworthy diagnostic tool that not only identifies the disease but also tells the farmer what to do about it.

This project builds a Crop Disease Detector, an AI-powered web application that lets a farmer upload a photo of a diseased leaf and instantly receive:
\begin{enumerate}
  \item The predicted disease (with confidence score).
  \item A visual explanation (Grad-CAM heatmap) showing which parts of the leaf the model focused on.
  \item Actionable treatment recommendations: specific fungicides, dosage, organic alternatives, and preventive measures tailored to the diagnosed disease.
\end{enumerate}

\section{Dataset}
\textbf{PlantVillage Dataset (Kaggle)}
\begin{itemize}
  \item $\sim$54,000 labeled leaf images.
  \item 38 disease classes across multiple crop species (tomato, potato, pepper, apple, grape, corn, etc.).
  \item Clean, well-documented, and widely used as a benchmark.
\end{itemize}
\textbf{Known limitation (which we address):} the dataset contains mostly lab-condition images with uniform backgrounds. Models trained on it can appear highly accurate while failing on real-world field photos. Our robustness testing (Section~\ref{sec:robust}) explicitly investigates this gap.

\section{Model and Architecture}
\textbf{Base model:} fine-tuned ResNet50 (or EfficientNet-B0) via transfer learning.

\textbf{Why transfer learning:}
\begin{itemize}
  \item Pretrained ImageNet weights provide strong low-level feature extractors (edges, textures, color gradients).
  \item Fine-tuning only the later layers substantially reduces training time and data requirements.
  \item Reported in the literature to reach high accuracy on PlantVillage with modest compute.
\end{itemize}

\textbf{Training pipeline:}
\begin{itemize}
  \item Input: $224\times224$ RGB images.
  \item Augmentation: random flips, rotations, color jitter, RandAugment.
  \item Loss: cross-entropy with class weights (to handle class imbalance).
  \item Optimizer: AdamW with cosine annealing.
  \item Early stopping on validation F1.
\end{itemize}

\section{Core Differentiator: Explainability and Treatment Engine}
Most student projects stop at ``predicted class + confidence.'' Ours goes two steps further.

\subsection{Grad-CAM Visual Explanations}
Grad-CAM generates a heatmap overlaid on the input image, highlighting the regions the model considered most important for its prediction.
\begin{itemize}
  \item \textbf{Builds trust:} the user can see whether the model is looking at the diseased spot rather than the background.
  \item \textbf{Supports bias analysis:} reveals whether the model learns dataset artifacts (e.g., uniform backgrounds) instead of disease patterns.
  \item \textbf{Debugging tool:} irrelevant highlighted regions signal the need for retraining or better augmentation.
\end{itemize}
\textit{Implementation:} \texttt{pytorch-grad-cam} or \texttt{captum}, integrated into the Streamlit UI.

\subsection{Treatment Recommendation Engine}
After diagnosis, the app displays a structured treatment plan for the detected disease, covering all 38 classes.

\begin{table}[h]
\centering
\begin{tabular}{@{}>{\bfseries}p{3.4cm}p{11cm}@{}}
\toprule
Field & Example (Early Blight -- Tomato) \\
\midrule
Cause & Fungus \textit{Alternaria solani} \\
Chemical treatment & Chlorothalonil or Mancozeb spray, every 7--10 days \\
Organic alternative & Neem oil spray, copper-based fungicide \\
Immediate action & Remove and destroy infected leaves; avoid overhead watering \\
Prevention & Crop rotation, resistant varieties, proper spacing for airflow \\
Severity guidance & If $>30\%$ of leaf area is affected, consider removing the plant \\
\bottomrule
\end{tabular}
\caption{Example treatment record. Final content will be verified against cited sources.}
\end{table}

\textbf{Implementation approach:}
\begin{itemize}
  \item \textbf{MVP:} curated JSON mapping each of the 38 classes to treatment info, sourced from agricultural extension documents (FAO, university guides).
  \item \textbf{Stretch:} RAG pipeline over agricultural PDFs using LangChain and a vector database for more context-aware advice.
\end{itemize}

\section{Real-World Robustness Testing}\label{sec:robust}
\textbf{The problem:} PlantVillage's clean lab images do not represent real field conditions: soil, shadows, multiple leaves, blur, and varying lighting.

\textbf{What we will do:}
\begin{enumerate}
  \item Collect $\sim$100 real-world field images (PlantDoc dataset, iNaturalist, or similar).
  \item Run the model on them and record the accuracy drop (hypothesis: a substantial drop relative to lab-condition accuracy).
  \item Apply mitigations: test-time augmentation, background removal (\texttt{rembg}), or fine-tuning on a small field-image set.
  \item Report before/after results.
\end{enumerate}

\section{User Interface (Streamlit)}
\textbf{Flow:}
\begin{enumerate}
  \item User opens the hosted web app.
  \item User uploads a leaf photo (drag-and-drop or camera capture on mobile).
  \item The app displays: predicted disease and confidence; Grad-CAM heatmap beside the original image; a structured treatment plan (chemical, organic, immediate, preventive); and a warning if confidence is below 70\%.
  \item Optional: download the diagnosis as a PDF report.
\end{enumerate}
\textbf{Hosting:} Streamlit Community Cloud (frontend) with the model hosted on Hugging Face Hub or GitHub Releases. No localhost-only evaluation.

\section{Evaluation and Bias Check}
\textbf{Metrics:} overall accuracy, macro F1, per-class precision/recall, $38\times38$ confusion matrix, top-3 accuracy.

\textbf{Bias check:}
\begin{itemize}
  \item Per-class performance breakdown to identify weak diseases.
  \item Class imbalance analysis: does the model favor majority classes?
  \item Grad-CAM analysis: does the model focus on disease regions or background artifacts?
  \item Robustness gap: lab vs.\ field image performance.
\end{itemize}

\section{Tech Stack}
\begin{table}[h]
\centering
\begin{tabular}{@{}ll@{}}
\toprule
Layer & Tool \\
\midrule
Model training & PyTorch, torchvision \\
Explainability & pytorch-grad-cam / captum \\
UI & Streamlit \\
Treatment data & Curated JSON (MVP) / LangChain + ChromaDB (stretch) \\
Deployment & Streamlit Community Cloud + Hugging Face \\
Version control & GitHub (datasets excluded via \texttt{.gitignore}) \\
Report & \LaTeX{} \\
\bottomrule
\end{tabular}
\end{table}

\section{Deliverables}
\begin{enumerate}
  \item GitHub repository: clean code, commit history showing individual contributions, README with setup instructions.
  \item Hosted web application with a live URL.
  \item \LaTeX{} technical report: methodology, architecture diagrams, Grad-CAM analysis, robustness findings, challenges faced.
  \item Treatment knowledge base (JSON or RAG pipeline), documented.
\end{enumerate}

\section{Team Roles}
\begin{table}[h]
\centering
\begin{tabular}{@{}>{\bfseries}p{3cm}p{11.4cm}@{}}
\toprule
Member & Primary responsibility \\
\midrule
Abdul Ahad (24L-3029) & Data pipeline, preprocessing, model training and evaluation \\
M. Anas (24L-3004) & Grad-CAM, treatment knowledge base, robustness testing \\
M. Babar Shahzad (24L-3047) & Streamlit UI, deployment, report assembly \\
\bottomrule
\end{tabular}
\caption{Initial split; all members review and understand the full codebase.}
\end{table}

\section{Timeline}
\begin{table}[h]
\centering
\begin{tabular}{@{}cp{9cm}c@{}}
\toprule
Step & Task & Days \\
\midrule
1 & Setup, dataset download, EDA & 1--2 \\
2 & Preprocessing, augmentation pipeline & 1--2 \\
3 & Baseline ResNet50 training & 2--3 \\
4 & Fine-tuning and hyperparameter search & 2--4 \\
5 & Grad-CAM integration & 1--2 \\
6 & Treatment engine (JSON + UI cards) & 1--2 \\
7 & Robustness testing on field images & 2--3 \\
8 & Streamlit UI and deployment & 1--2 \\
9 & Evaluation, bias check, \LaTeX{} report & 2--3 \\
-- & Buffer for debugging & 2--3 \\
\bottomrule
\end{tabular}
\caption{Tentative schedule; adjusted to the instructor's deliverable dates.}
\end{table}

\section{Risks and Mitigations}
\begin{table}[h]
\centering
\begin{tabular}{@{}p{6cm}p{9.4cm}@{}}
\toprule
Risk & Mitigation \\
\midrule
Model overfits to clean backgrounds & Aggressive augmentation and robustness testing \\
Class imbalance skews results & Weighted loss and per-class metrics \\
Deployment issues & Streamlit Cloud (low-config) with model on Hugging Face \\
Treatment information accuracy & Source from FAO / university extension documents; cite sources \\
Time overrun & MVP first (classifier + Grad-CAM + JSON treatments), then stretch goals \\
\bottomrule
\end{tabular}
\end{table}

\end{document}
